\documentclass[11pt]{article}

\usepackage[margin=1in]{geometry}
\usepackage{amsmath,amssymb,amsthm}
\usepackage{hyperref}

\title{The RSK Transformer as a Manifold Microscope}
\author{}
\date{}

\begin{document}
\maketitle

\section*{The claim}

A 1.2M-parameter transformer\footnote{Model and checkpoints available at \url{https://huggingface.co/RobBobin/rsk-transformer}.} trained on inverse RSK is an unusually
good instrument for studying how neural networks learn maps between
manifolds. Three properties make it exceptional: the input and output
spaces have known combinatorial structure, the model is small enough
for complete mechanistic analysis, and the same architecture learns
four distinct bijections that share a common local geometry.

\section*{Why ``manifold'' is the right word}

Standard Young Tableaux are discrete objects, but the structured
embedding---value, row, column, tableau identity, each mapped to
$\mathbb{R}^{128}$ and summed---lifts them onto a continuous
submanifold $M_{\mathrm{in}} \subset \mathbb{R}^{128 \times 2n}$. The
SYT conditions (rows increase left-to-right, columns increase
top-to-bottom) constrain which points on this manifold are occupied.
The transformer's six residual layers then implement a smooth flow
\[
  x \mapsto x + f(x)
\]
that deforms $M_{\mathrm{in}}$ into a manifold $M_{\mathrm{out}}$
where the target permutation is linearly separable. This is not
metaphor: each layer is a composition of differentiable operations
(linear projections, softmax, GELU, layer norm), so the composite map
is smooth, and the residual connection makes each layer a
near-identity diffeomorphism.

\section*{What makes it a microscope}

Most transformer interpretability work studies language models with
billions of parameters, polysemantic neurons, and no ground truth
about what the ``right'' algorithm is. The RSK transformer inverts all
three obstacles.

\paragraph{Known ground truth.}
The inverse RSK algorithm is a well-defined sequence of reverse
row-insertions, fully specified by Schensted's 1961 construction. When
we find SAE features that fire on specific tableau entries ($F_{356}$
on $P[7]$, $F_{622}$ on $Q[0]$) or on Burge local-rule triples in the
cylindric model, we can check these against the textbook algorithm.
The manifold has a known atlas.

\paragraph{Complete ablation is feasible.}
With 6 layers, 8 heads, and 1.2M parameters, every component can be
individually ablated, and every attention pattern can be visualised.
The finding that Layer~2 is critical (ablation drops accuracy from
100\% to 62\%) while Layers~4--5 are vestigial tells us exactly where
the flow does its work. In a billion-parameter model, this kind of
localisation is usually impossible.

\paragraph{The anti-memorisation argument is clean.}
At $n = 10$, the model trains on 500K of 3.6M possible permutations
(14\%) and achieves 100\% on the held-out 86\%. At $n = 15$, it trains
on 0.00004\% of the $1.3 \times 10^{12}$ permutation space and still
reaches 99.99\%. A lookup table would need orders of magnitude more
parameters. The model must have learned the manifold's
geometry---the smooth structure that allows interpolation between
training points---rather than memorising input--output pairs.

\section*{The fiber bundle decomposition}

The architecture has a natural geometric interpretation as a map on a
trivial fiber bundle $B \times \mathbb{R}^{128}$, where $B$ is the
set of $2n$ token positions and $\mathbb{R}^{128}$ is the embedding
fiber at each position.

Attention is the horizontal component: it transports information
between positions, implementing parallel transport along paths
determined by the data. The MLP is the vertical component: it
transforms within each fiber independently. This is not a loose
analogy---it is the standard decomposition of a map on a fiber bundle
into horizontal and vertical parts, and the architecture enforces it
by construction.

The SAE results on the cylindric model give this decomposition teeth.
Features fire on the specific triples of positions where the Burge
local rule operates, at $3.0$--$3.3\times$ above baseline lift. The
connection on the bundle---the pattern of which positions attend to
which---is not arbitrary. It follows the combinatorial paths that the
growth diagram algorithm prescribes.

\section*{The attention entropy phase transition}

At $n = 8$ and $n = 10$, there is a sharp boundary at
Layer~$3 \to$~Layer~4: early layers have concentrated (low-entropy)
attention, late layers have diffuse (high-entropy) attention. In the
flow picture, this is a fixed point: the manifold has been deformed
into its target shape, and subsequent layers produce negligible
displacement.

At $n = 15$, the phase transition disappears. All six layers maintain
structured attention. The manifold is more complex, the flow needs
more time to converge, so it extends to fill the available depth. This
is the continuous-flow analogue of a stiff ODE requiring smaller time
steps, and it is direct evidence that the model allocates computation
adaptively---using exactly as many ``structured layers'' as the
problem demands.

\section*{Four bijections, one local geometry}

The strongest evidence for the manifold perspective is that the same
architecture, with no modification, learns four different bijections:
\begin{enumerate}
  \item Permutation RSK ($\sigma \leftrightarrow (P, Q)$ standard Young tableaux)
  \item Matrix RSK (integer matrix $\leftrightarrow (P, Q)$ semistandard Young tableaux)
  \item Hillman--Grassl (reverse plane partition $\leftrightarrow$ arbitrary filling)
  \item Cylindric growth diagrams (cylindric plane partition $\leftrightarrow$ ALCD labels)
\end{enumerate}

These bijections act on different combinatorial objects with different
symmetries. But they all share the same local structure: Fomin's
growth diagrams, in which a grid is filled vertex-by-vertex using a
local rule that takes three neighbouring partitions and produces the
fourth.

In manifold terms: the input manifolds are different, but they have
the same local differential structure. The structured 2D embedding
preserves this local geometry, and the transformer learns a flow that
works for any manifold with Fomin-type local structure. The
architecture is a microscope focused on exactly the right geometric
feature.

\section*{What this suggests}

Small models trained on well-understood mathematical bijections may be
the cleanest setting we have for studying how neural networks learn
maps between structured spaces. The RSK transformer offers something
rare: a model where we know what the manifolds are, we know what the
correct map between them is, we can verify that the model has learned
it, and we can trace how each layer deforms the geometry. If we want
to understand what transformers do to manifolds, we should start here.

\end{document}
